

%%%% MATH %%%%
\usepackage{bm} % bold symbols using \bm
\renewcommand{\vec}[1]{\mathbf{#1}}
\newcommand{\rmi}{\mathrm{i}}
\newcommand{\rme}{\mathrm{e}}
\newcommand{\rmd}{\,\mathrm{d}}
% \DeclareMathOperator{\rmd}{d}%gives wrong spacing after \rmd...
\newcommand{\rmb}{\mathrm{b}}
\newcommand{\rmf}{\mathrm{f}}
\newcommand{\dg}{\ensuremath{^\circ}}
\newcommand{\uvec}[1]{\vec{\hat{#1}}}

\usepackage[]{amsmath}
%\newcommand\tline[2]{$\underset{\text{#1}}{\text{\underline{\hspace{#2}}}}$}
\newcommand\tline[2]{$\underset{\text{#2}}{\text{\underline{#1}}}$}

%\renewcommand{\sqrt}[1]{\ensuremath (#1)^{1/2}}
%\renewcommand{\sqrt}[1]{\ensuremath \textrm{((}#1\textrm{))}^{1/2}}
%\renewcommand{\sqrt}[1]{\ensuremath \Big(#1 \Big)^{1/2}}
\renewcommand{\sqrt}[1]{\ensuremath \Big(#1 \Big)^{1/2}}
\newcommand{\Sqrt}[1]{\ensuremath \bigg(#1 \bigg)^{1/2}}

\renewcommand{\langle}{\ensuremath <}
\renewcommand{\rangle}{\ensuremath >}
%\tline{Repetition is needed}{$F = ma$}

% https://tex.stackexchange.com/questions/22773/making-a-big-summation-sign
%\usepackage{calc}
\newlength{\depthofsumsign}
\setlength{\depthofsumsign}{\depthof{$\sum$}}
\newlength{\totalheightofsumsign}
\newlength{\heightanddepthofargument}

\newcommand{\nsum}[1][1.4]{
  \mathop{%
    \raisebox
    {-#1\depthofsumsign+1.4\depthofsumsign}%-#1\depthofsumsign+1\depthofsumsign
    {\scalebox
      {#1}
      {$\displaystyle\sum$}%
    }
  }
}

\newcommand{\nint}[1][1.4]{
  \mathop{%
    \raisebox
    {-#1\depthofsumsign+1\depthofsumsign}
    {\scalebox
      {#1}
      {$\displaystyle\int$}%
    }
  }
}

\newcommand{\niint}[1][1.4]{
  \mathop{%
    \raisebox
    {-#1\depthofsumsign+1\depthofsumsign}
    {\scalebox
      {#1}
      {$\displaystyle\iint$}%
    }
  }
}

%\newlength{\depthofparenthesis}
%\setlength{\depthofparenthesis}{\depthof{$($}}
%\newcommand{\vasen}[1][1.4]{
%  \mathop{%
%    \raisebox
%    {-#1\depthofparenthesis+1\depthofparenthesis}
%    {\scalebox
%      {#1}
%      {$\displaystyle($}%
%    }
%  }
%}

%% enabling use of \bigg( and \bigg)^{3/2} etc.
\renewcommand\big[1]{\raisebox{-1.33\depthofsumsign+1\depthofsumsign}{\scalebox{1.333}{$\displaystyle \text{\bigsymbolsfont#1}$}}%
}
\renewcommand\Big[1]{\raisebox{-1.75\depthofsumsign+1\depthofsumsign}{\scalebox{1.75}{$\displaystyle \text{\bigsymbolsfont#1}$}}%
}
\renewcommand\bigg[1]{\raisebox{-2.2\depthofsumsign+1\depthofsumsign}{\scalebox{2.2}{$\displaystyle \text{\bigsymbolsfont#1}$}}%
}
\renewcommand\Bigg[1]{\raisebox{-2.7\depthofsumsign+1\depthofsumsign}{\scalebox{2.7}{$\displaystyle \text{\bigsymbolsfont#1}$}}%
}

\renewcommand\bigl[1]{\raisebox{-1.33\depthofsumsign+1\depthofsumsign}{\scalebox{1.333}{$\displaystyle \text{\bigsymbolsfont#1}$}}%
}
\renewcommand\Bigl[1]{\raisebox{-1.75\depthofsumsign+1\depthofsumsign}{\scalebox{1.75}{$\displaystyle \text{\bigsymbolsfont#1}$}}%
}
\renewcommand\biggl[1]{\raisebox{-2.2\depthofsumsign+1\depthofsumsign}{\scalebox{2.2}{$\displaystyle \text{\bigsymbolsfont#1}$}}%
}
\renewcommand\Biggl[1]{\raisebox{-2.7\depthofsumsign+1\depthofsumsign}{\scalebox{2.7}{$\displaystyle \text{\bigsymbolsfont#1}$}}%
}

\renewcommand\bigr[1]{\raisebox{-1.33\depthofsumsign+1\depthofsumsign}{\scalebox{1.333}{$\displaystyle \text{\bigsymbolsfont#1}$}}%
}
\renewcommand\Bigr[1]{\raisebox{-1.75\depthofsumsign+1\depthofsumsign}{\scalebox{1.75}{$\displaystyle \text{\bigsymbolsfont#1}$}}%
}
\renewcommand\biggr[1]{\raisebox{-2.2\depthofsumsign+1\depthofsumsign}{\scalebox{2.2}{$\displaystyle \text{\bigsymbolsfont#1}$}}%
}
\renewcommand\Biggr[1]{\raisebox{-2.7\depthofsumsign+1\depthofsumsign}{\scalebox{2.7}{$\displaystyle \text{\bigsymbolsfont#1}$}}%
}


% bigger symbols, e.g. ( ) , [ ] , { , } etc.
\newcommand{\bg}[2][1.4]{
  \mathop{%
    \raisebox
    {-#1\depthofsumsign+1\depthofsumsign}
    {\scalebox
      {#1}
      {$\displaystyle \appear{#2} $}%
    }
  }
}

%\usepackage{mathtools} 
%\DeclarePairedDelimiter\autobracket{(}{)}
%\newcommand{\br}[1]{\autobracket*{#1}}
%\abs[\Bigg]{\frac{a}{b}}

% \newcommand{\T}{\intercal}

\newcommand{\ket}[1]{$\left|#1\right\rangle$}

\newcommand\varpm{\mathbin{\vcenter{\hbox{%
  \oalign{\hfil$\scriptstyle+$\hfil\cr
          \noalign{\kern-.3ex}
          $\scriptscriptstyle({-})$\cr}%
}}}}
\newcommand\varmp{\mathbin{\vcenter{\hbox{%
  \oalign{$\scriptstyle({+})$\cr
          \noalign{\kern-.3ex}
          \hfil$\scriptscriptstyle-$\hfil\cr}%
}}}}

%% cardinal sine seems to be missing a definition..
%% the def should make a small space after the sinc!!!
% \newcommand{\sinc}{\mathrm{sinc}}
% \newcommand{\sinc}{\,\mathrm{sinc}}
\DeclareMathOperator{\sinc}{sinc}
%%

%%
\newcommand{\tabtwoline}[2]{
\begin{tabular}{@{}c@{}} #1 \\ #2 \end{tabular}
}
\newcommand{\tabthreeline}[3]{
\begin{tabular}{@{}c@{}c@{}} #1 \\ #2 \\ #3 \end{tabular}
}
\newcommand{\tabfourline}[4]{
\begin{tabular}{@{}c@{}c@{}c@{}} #1 \\ #2 \\ #3 \\ #4 \end{tabular}
}

\usepackage{pgffor}

%% partial derivatives:
\newcommand{\pder}[2][]{\frac{\partial #1}{\partial #2}}
\newcommand{\ppder}[2][]{\frac{\partial^2 #1}{\partial #2^2}}

%% total derivatives:
\usepackage{commath}
\newcommand{\der}[2][]{\frac{\rmd #1}{\rmd #2}}
\newcommand{\dder}[2][]{\frac{\rmd^2 #1}{\rmd #2^2}}

%% macros for vectors:
% https://tex.stackexchange.com/questions/270223/macro-for-writing-vectors
\usepackage{stackengine}
\newcommand\Vector[1]{\setstackEOL{,}\bracketVectorstack{#1}}

%%%%%%%%%%%%%%%%%%%% MUELLER MATRICES %%%%%%%%%%%%%%%%%%%%
\newcommand{\MuellerHlin}{
    $\frac{1}{2} 
    \begin{bmatrix}
        1 & 1 & 0 & 0 \\
        1 & 1 & 0 & 0 \\
        0 & 0 & 0 & 0 \\
        0 & 0 & 0 & 0 \\
    \end{bmatrix}
    $
}
\newcommand{\MuellerVlin}{
    $\frac{1}{2} 
    \begin{bmatrix}
        1 & -1 & 0 & 0 \\
        -1 & 1 & 0 & 0 \\
        0 & 0 & 0 & 0 \\
        0 & 0 & 0 & 0 \\
    \end{bmatrix}
    $
}
\newcommand{\Muellerdiaglin}{
    $\frac{1}{2} 
    \begin{bmatrix}
        1 & 0 & 1 & 0 \\
        0 & 0 & 0 & 0 \\
        1 & 0 & 1 & 0 \\
        0 & 0 & 0 & 0 \\
    \end{bmatrix}
    $
}
\newcommand{\Muellermdiaglin}{
    $\frac{1}{2} 
    \begin{bmatrix}
        1 & 0 & -1 & 0 \\
        0 & 0 & 0 & 0 \\
        -1 & 0 & 1 & 0 \\
        0 & 0 & 0 & 0 \\
    \end{bmatrix}
    $
}
\newcommand{\MuellerQWPHlin}{
    $\begin{bmatrix}
        1 & 0 & 0 & 0 \\
        0 & 1 & 0 & 0 \\
        0 & 0 & 0 & 1 \\
        0 & 0 & -1 & 0 \\
    \end{bmatrix}
    $
}
\newcommand{\MuellerQWPVlin}{
    $\begin{bmatrix}
        1 & 0 & 0 & 0 \\
        0 & 1 & 0 & 0 \\
        0 & 0 & 0 & -1 \\
        0 & 0 & 1 & 0 \\
    \end{bmatrix}
    $
}

%%%%%%%%%%%%%%%%%%%% STOKES MATRICES %%%%%%%%%%%%%%%%%%%%
\newcommand{\StokesHlin}{
    $\begin{bmatrix}
        1 & 0   \\
        0 & 0   \\
    \end{bmatrix}
    $
}

\newcommand{\StokesYlin}{
    $\begin{bmatrix}
        0 & 0   \\
        0 & 1   \\
    \end{bmatrix}
    $
}

\newcommand{\Stokesdiaglin}{
    $\frac{1}{2} 
    \begin{bmatrix}
        1 & 1   \\
        1 & 1   \\
    \end{bmatrix}
    $
}
\newcommand{\Stokesmdiaglin}{
    $\frac{1}{2} 
    \begin{bmatrix}
        1 & -1   \\
        -1 & 1   \\
    \end{bmatrix}
    $
}

\newcommand{\StokesQWPVlin}{
    $\rme^{\rmi \pi / 4} 
    \begin{bmatrix}
        1 & 0   \\
        0 & -\rmi   \\
    \end{bmatrix}
    $
}

\newcommand{\StokesQWPHlin}{
    $\rme^{\rmi \pi / 4} 
    \begin{bmatrix}
        1 & 0   \\
        0 & \rmi   \\
    \end{bmatrix}
    $
}

%%%%%%%%%%%%%%%%%%%% STOKES VECTORS %%%%%%%%%%%%%%%%%%%%
\newcommand{\Stokesvector}[4]{
    $\begin{bmatrix}
        #1   \\
        #2   \\
        #3   \\
        #4   \\
    \end{bmatrix}
    $
}

%%%%%%%%%%%%%%%%%%%% JONES VECTORS %%%%%%%%%%%%%%%%%%%%
\newcommand{\Jonesvector}[3]{
    $#1 \begin{bmatrix}
        #2   \\
        #3   \\
    \end{bmatrix}
    $
}
\usepackage{booktabs}

